\section{数据与测度}

\subsection{空气质量面板}\label{sec:panel}

数据为全国城市空气质量逐小时监测，覆盖 \textbf{389 个城市}、\textbf{2015—2026 年 12 个春节窗口}。指标包括细颗粒物 \pmm{2.5}、可吸入颗粒物 \pmm{10}、空气质量指数（AQI）、\SOtwo{}、\NOtwo{}、\Othree{} 与 \COone{}。

每个春节年份展开为一个\textbf{相对日}窗口：正月初一记为 \texttt{rel\_day} $=+1$，覆盖 $-7$ 至 $+7$。以初一而非除夕为基准日，是因为实测峰值出现在\textbf{正月初一}，不在除夕夜。这与“跨年燃放”的直觉略有出入，但对测度很重要：若以除夕为峰，窗口内的高点会被系统性错置一天。

\textbf{基线}定义为除夕前 7 日至前 2 日的日均值（\texttt{rel\_day} $=-7$ 至 $-2$）。\textbf{主结果变量}为初一日均 \pmm{2.5} 与\textbf{尖峰增量}（初一减基线）。

本文使用三个样本口径，正文中始终标明：

\begin{table}[H]
\centering
\caption{三个样本口径与各自用途}
\label{tab:samples}
\small
\begin{tabular}{L{3.6cm}L{3.2cm}L{3.0cm}L{4.2cm}}
\toprule
\textbf{口径} & \textbf{城市数} & \textbf{城市-年数} & \textbf{用途} \\
\midrule
平衡面板 & 306 & $306\times12=3{,}672$ & 估计与方差分解的主口径；判据是初一与基线日均在全部年份俱全 \\
非平衡面板 & 389 & 4{,}240 & 描述性统计的补充口径；判据是凡初一与基线俱全的城市-年均计入 \\
作图画板 & 308 & — & 仅用于图形；判据是“该年可用天数 $\geq 13$” \\
\bottomrule
\end{tabular}
\end{table}

\noindent 三个口径互不相同，混用会导致数字对不上。本文的规则是：\textbf{凡与估计有关的量用平衡面板（3,672）；描述“全国城市-年的分布”时，同时报告非平衡口径（4,240）并注明}。第~\ref{sec:magnitude}~节的两组数字按此规则给出。全样本城市-日共 63,515 行，属于 389 城非平衡面板（4,240 个城-年 $\times$ 15 日），不是作图画板。

\subsection{政策库}

\texttt{policy/policy.csv} 一行记录一次政策变动。字段包括政策类型（无专门管制／限放／禁放）、生效日期、地理范围（全市域／多个区县／仅中心城区）、罚则、发文主体和是否使用税收工具。\textbf{证据分档是硬性纪律}：只有核过一手来源的行才标 \texttt{verified=Y}；媒体摘要一律标 \texttt{N}，并注明来源性质。

这一纪律不是形式主义。政策库中曾出现三类错误：把“专项”人数当成“总量”的子集；数字正确但来源链接指向错误年份；来源链接失效且无存档。\textbf{三类错误只能靠回原文发现，不能靠核对数字发现。}

\subsection{政策文件的全文取证}\label{sec:fulltext}

对 34 个城市的大气污染防治方案逐份下载原文，归一化后检索“烟花爆竹”，记录命中位置、是否为独立任务、牵头单位。\textbf{“未找到”与“没有”严格区分}：官方全文零命中记为“确证没有”；文件未能获取的记为“未找到”，不计入分母。这一区分使 34 城的比例成为可辩护的数字，不是检索能力的函数。逐城结果见附录~\ref{app:documents}。

\subsection{执法成本数据}

执法数据的采集结果与预期相反。原设计假定“官方通报几乎年年发”，可以构造逐城逐年的成本面板。实测三批、覆盖 21 个城市后，\texttt{cases} 字段只有 28 格有值，其中仅 8 格是“含违规燃放”的公安查处案件数，其余案由各异。山东 6 市呈\textbf{结构性稀疏}：淄博六年中四年市级通稿通篇无“烟花爆竹”四字；烟台六年市级通报从未给出“查处 X 起”；滨州 2024 年（全年禁放首年）市级通稿完全没有烟花爆竹内容。

\textbf{结论是：成本侧不能做逐城逐年面板。} 它能支撑存在性证据与少数城市的个案量级。这一结论改写了原设计，本身也是本文的一个发现：\textbf{这项政策的执法成本很高，却系统性缺乏核算。}
